% 00-resumo.tex -- incluído por main-pt.tex
\begin{abstract-pt}
Registros de créditos de carbono baseados em blockchain vêm sendo propostos para conferir integridade e rastreabilidade a projetos REDD+, e a Lei nº 15.042/2024 cria a perspectiva de um registro nacional no Brasil. A viabilidade do consenso bizantino nessa escala, contudo, tem sido discutida com modelos que dimensionam o número de validadores pela área florestal e que estimam o tráfego de consenso em forma fechada. Este trabalho investiga em que condições cada família de protocolos de consenso permanece viável em um registro nacional e qual restrição se esgota primeiro. O número de validadores é tratado como uma escala de governança, separada da escala de mercado que determina a demanda, e a demanda é modelada a partir da estrutura de conciliação periódica do mercado regulado. Protocolos admissíveis são selecionados por requisitos do domínio e avaliados no simulador JABS, com tráfego medido em nível de pacote e validadores bizantinos, em um experimento fatorial cujo protocolo de análise é fixado antes da bateria principal. Elasticidades em relação às duas escalas são estimadas por metamodelos validados fora da faixa de ajuste, e as interações são testadas por ANOVA com verificações de robustez. \tbd{Elasticidades medidas, restrição vinculante por regime e fronteiras de viabilidade.} \tbd{Confirmar ou refutar a hipótese de que a viabilidade depende mais do desenho de governança do que do volume de transações.} Em qualquer caso, as fronteiras obtidas indicam quantos validadores um registro pode comportar para cada nível de demanda.
\end{abstract-pt}

\begin{abstract-en}
Blockchain-based carbon credit registries have been proposed to provide integrity and traceability to REDD+ projects, and Brazilian Law 15,042/2024 opens the prospect of a national registry. The viability of Byzantine consensus at this scale, however, has been discussed with models that size the validator set by forest area and estimate consensus traffic in closed form. This work investigates under which conditions each family of consensus protocols remains viable for a national registry and which constraint becomes binding first. The number of validators is treated as a governance scale, separated from the market scale that determines demand, and demand is modeled from the periodic reconciliation structure of the regulated market. Admissible protocols are selected from domain requirements and evaluated on the JABS simulator, with packet-level traffic measurement and Byzantine validators, in a factorial experiment whose analysis protocol is fixed before the main battery. Elasticities with respect to both scales are estimated by metamodels validated outside the fitting range, and interactions are tested by ANOVA with robustness checks. \tbd{Measured elasticities, binding constraint per regime and feasibility frontiers.} \tbd{Confirm or refute the hypothesis that viability depends more on governance design than on transaction volume.} In either case, the resulting frontiers indicate how many validators a registry can accommodate for each demand level.
\end{abstract-en}

\begin{pchaves}
Blockchain, Consenso Bizantino, Escalabilidade, REDD+, Registro de Carbono, Análise de Variância
\end{pchaves}

\begin{keywords}
Blockchain, Byzantine Consensus, Scalability, REDD+, Carbon Registry, Analysis of Variance
\end{keywords}
